\section{Hypothesis Testing and Root Cause}

\subsection{Candidate Causes}
\begin{table}[H]
    \centering
    \small
    \begin{tabular}{p{3.3cm} p{2.1cm} p{2.2cm} p{3.9cm}}
        \toprule
        \textbf{Hypothesis} & \textbf{Likelihood} & \textbf{Evidence Status} & \textbf{Comment} \\
        \midrule
        Control-loop instability during load transition & Medium & Open & Requires replay under instrumented conditions. \\
        Component overstress on switching node & High & Supported & Correlates with hotspot and event timing. \\
        Fixture or cabling artifact & Low & Partially ruled out & Reproduction with alternate setup still recommended. \\
        Firmware fault handling defect & Medium & Open & Needs trace review and fault-code decode. \\
        \bottomrule
    \end{tabular}
    \caption{Root-cause hypothesis table}
    \label{tab:root-cause-hypotheses}
\end{table}

\subsection{Structured Reasoning}
\begin{enumerate}
    \item State the strongest evidence supporting each hypothesis.
    \item State the evidence that contradicts it.
    \item Define the next experiment that would separate the leading candidates.
    \item Avoid collapsing multiple contributing causes into a single vague label.
\end{enumerate}

\decisionbox{
If a root cause is claimed, name the direct cause, contributing condition, and
missing control separately. This avoids weak conclusions such as ``component
failed'' without explaining why.
}
